% Part: turing-machines
% Chapter: undecidability
% Section: verification

\documentclass[../../../include/open-logic-section]{subfiles}

\begin{document}

\olfileid[hi]{tur}{und}{ver}
\olsection{निरूपण का सत्यापन}

\begin{explain}
यह सत्यापित करने के लिए कि हमारा निरूपण काम करता है, हमें दो बातें
सिद्ध करनी हैं। पहले हमें दिखाना है कि यदि $M$ आगत~$w$ पर रुकती है, तो
$!T(M, w) \lif !E(M, w)$ वैध है। फिर हमें इसका विपरीत सिद्ध करना है,
अर्थात् यदि $!T(M, w) \lif !E(M, w)$ वैध है, तो आगत~$w$ पर चलाने पर
$M$ वास्तव में अंततः रुकती है।

इन दोनों बातों को सिद्ध करने की रणनीतियाँ बहुत भिन्न हैं। पहले परिणाम
के लिए हमें दिखाना है कि प्रथम-कोटि तर्कशास्त्र का !!a{sentence}
(अर्थात् $!T(M, w) \lif !E(M, w)$) वैध है। ऐसा करने की सबसे सरल रीति
!!a{derivation} देना है। किंतु हमारे प्रमाण को सभी $M$ और $w$ के लिए
काम करना है, इसलिए वास्तव में कोई एक !!{sentence} नहीं है जिसके लिए
हमें !!a{derivation} देना हो, बल्कि ऐसे वाक्य अनंततः अनेक हैं।
अतः अधिकतम इतना किया जा सकता है कि आगमन से सिद्ध करें: $M$ और~$w$ चाहे
जैसे हों और आगत~$w$ पर $M$ को रुकने में चाहे जितने चरण लगें,
$!T(M, w) \lif !E(M, w)$ की !!a{derivation} होगी।

स्वाभाविक है कि हमारा आगमन उन चरणों की संख्या पर चलेगा जो $M$ किसी
रुकने वाले विन्यास तक पहुँचने से पहले लेती है। आगमनात्मक प्रमाण में हम
स्थापित करेंगे कि आगत~$w$ पर $M$ के निष्पादन के प्रत्येक चरण~$n$ के लिए
$!T(M, w) \Entails !C(M, w, n)$, जहाँ $!C(M, w, n)$ आगत~$w$ पर चलाई
गई $M$ के $n$ चरणों के बाद वाले विन्यास का ठीक वर्णन करता है। अब यदि
$M$ आगत~$w$ पर, मान लें, $n$ चरणों के बाद रुकती है, तो $!C(M, w, n)$
रुकने वाले विन्यास का वर्णन करेगा। हम यह भी दिखाएँगे कि जब भी
$!C(M, w, n)$ रुकने वाले विन्यास का वर्णन करता है, तब
$!C(M, w, n) \Entails !E(M, w)$। अतः यदि $M$ आगत~$w$ पर रुकती है, तो
किसी~$n$ के लिए $M$, $n$ चरणों के बाद रुकने वाले विन्यास में होगी। फलतः
$!T(M, w) \Entails !C(M, w, n)$, जहाँ $!C(M, w, n)$ रुकने वाले विन्यास
का वर्णन करता है; और चूँकि उस दशा में
$!C(M, w, n) \Entails !E(M, w)$, हमें $T(M, w) \Entails !E(M, w)$
मिलता है, अर्थात् $\Entails !T(M, w) \lif !E(M, w)$।

विपरीत दिशा की रणनीति बहुत भिन्न है। यहाँ हम मानते हैं कि
$\Entails !T(M, w) \lif !E(M, w)$ और हमें सिद्ध करना है कि $M$
आगत~$w$ पर रुकती है। परिकल्पना से मिलता है कि
$!T(M, w) \Entails !E(M, w)$, अर्थात् हर उस !!{structure} में
$!E(M, w)$ सत्य है जिसमें $!T(M, w)$ सत्य है। इसलिए हम
!!a{structure}~$\Struct{M}$ का वर्णन करेंगे जिसमें $!T(M, w)$ सत्य है:
उसका प्रांत $\Nat$ होगा, और सभी $\Obj Q_q$ तथा $\Obj S_\sigma$ की
व्याख्या आगत~$w$ पर $M$ के निष्पादन के दौरान आने वाले विन्यासों द्वारा
दी जाएगी। उदाहरणार्थ, $\Sat{M}{\Obj Q_q(\num{m}, \num{n})}$ तभी जब
आगत~$w$ पर $n$ चरणों तक चलाए जाने पर $T$, अवस्था~$q$ में हो और
खाना~$m$ पढ़ रही हो। अब परिकल्पना के अनुसार
$!T(M, w) \Entails !E(M, w)$, और रचना के अनुसार
$\Sat{M}{!T(M, w)}$, इसलिए $\Sat{M}{!E(M, w)}$। किंतु
$\Sat{M}{!E(M, w)}$ तभी जब कोई $n \in \Domain{M} = \Nat$ ऐसा हो कि
आगत~$w$ पर चलाई गई $M$, $n$ चरणों के बाद रुकने वाले विन्यास में हो।
\end{explain}

\begin{defn}
$!C(M, w, n)$ को निम्न !!{sentence} मानें:
\[
\Obj Q_q(\num{m}, \num{n}) \land \Obj S_{\sigma_0}(\num{0}, \num{n})
\land \dots \land \Obj S_{\sigma_k}(\num{k}, \num{n}) \land
\lforall[x][(\num{k} < x \lif \Obj S_\TMblank(x, \num{n}))]
\]
जहाँ समय~$n$ पर $M$ की अवस्था $q$ है, समय~$n$ पर $M$ खाना~$m$ पढ़ रही
है, समय~$n$ पर खाना~$i$ में प्रतीक~$\sigma_i$ है ($0 \le i \le k$),
और $k$ समय~$0$ पर पट्टी का सबसे दायाँ अरिक्त खाना अथवा $n$ चरणों के बाद
तक पट्टी-शीर्ष द्वारा देखा गया सबसे दायाँ खाना है---इन दोनों में जो भी
अधिक दायाँ हो।
\end{defn}

\begin{lem}\ollabel{lem:halt-config-implies-halt}
यदि आगत~$w$ पर चलाई गई $M$, $n$ चरणों के बाद रुकने वाले विन्यास में है,
तो $!C(M, w, n) \Entails !E(M, w)$।
\end{lem}

\begin{proof}
मान लें कि $M$ आगत~$w$ पर $n$ चरणों के बाद रुकती है। कोई अवस्था~$q$,
खाना~$m$, और प्रतीक~$\sigma$ ऐसे हैं कि:
\begin{enumerate}
\item $n$ चरणों के बाद $M$, अवस्था~$q$ में है और खाना~$m$ पढ़ रही है,
  जिस पर~$\sigma$ लिखा है।
\item संक्रमण फलन $\delta(q, \sigma)$ अपरिभाषित है।
\end{enumerate}
$!C(M, w, n)$ इस विन्यास का वर्णन है और इसमें खंड
$\Obj Q_{q}(\num{m}, \num{n})$ तथा
$\Obj S_{\sigma}(\num{m}, \num{n})$ सम्मिलित होंगे। ये खंड मिलकर
$!E(M, w)$ को फलित करते हैं:
\[
\lexists[x][\lexists[y][(\bigvee_{\tuple{q, \sigma} \in
      X}(\Obj Q_q(x, y) \land \Obj S_\sigma(x, y)))]]
\]
क्योंकि $\tuple{q',\sigma'} \in X$, अतः
$\Obj Q_{q'}(\num{m}, \num{n}) \land S_{\sigma'}(\num{m},
\num{n}) \Entails \bigvee_{\tuple{q, \sigma} \in X}
(\Obj Q_q(\num{m}, \num{n}) \land \Obj S_{\sigma}(\num{m}, \num{n}))$।
\end{proof}

\begin{explain}
इसलिए यदि $M$ आगत $w$ पर रुकती है, तो कोई~$n$ ऐसा है कि
$!C(M, w, n) \Entails !E(M,w)$। अब हम दिखाएँगे कि किसी भी समय~$n$ के
लिए $!T(M, w) \Entails !C(M, w, n)$।
\end{explain}

\begin{lem}
\ollabel{lem:config}
प्रत्येक $n$ के लिए, यदि $M$, $n$ चरणों के बाद नहीं रुकी है, तो
$!T(M, w) \Entails !C(M, w, n)$।
\end{lem}

\begin{proof}
आगमन का आधार: यदि $n = 0$, तो $!C(M, w, 0)$ के संयोजक
$!T(M, w)$ के भी संयोजक हैं और इसलिए उससे फलित होते हैं।

आगमन-परिकल्पना: यदि $M$, $n$वें चरण से पहले नहीं रुकी है, तो
$!T(M,w) \Entails !C(M, w, n)$। हमें दिखाना है कि (जब तक
$!C(M, w, n)$ किसी रुकने वाले विन्यास का वर्णन न करे)
$!T(M, w) \Entails !C(M, w, n+1)$।

मान लें कि $n > 0$ और $w$ पर आरंभ की गई $M$, $n$ चरणों के बाद
अवस्था~$q$ में है और खाना~$m$ पढ़ रही है। चूँकि $M$, $n$ चरणों के बाद
नहीं रुकती, $M$ के प्रोग्राम में निम्न तीन रूपों में से किसी एक का निर्देश
अवश्य होना चाहिए:

\begin{enumerate}
\item \ollabel{right} $\delta(q, \sigma) = \tuple{q', \sigma', \TMright}$

\item \ollabel{left} $\delta(q, \sigma) = \tuple{q', \sigma', \TMleft}$

\item \ollabel{stay} $\delta(q, \sigma) = \tuple{q', \sigma', \TMstay}$
\end{enumerate}

अब हम इन तीनों स्थितियों पर बारी-बारी से विचार करेंगे।

\begin{enumerate}
\item मान लें कि \olref{right} रूप का कोई निर्देश है।
  \olref[rep]{defn:tm-descr}\olref[rep]{rep-right} के अनुसार इसका अर्थ
  है कि
\begin{align*}
& \lforall[x][\lforall[y][((\Obj Q_{q}(x,
  y) \land \Obj S_\sigma(x, y)) \lif {}]]\\
& \qquad (\Obj Q_{q'}(x',
  y') \land \Obj S_{\sigma'}(x, y') \land !A(x, y)))
\intertext{$!T(M,w)$ का एक संयोजक है। इससे निम्न !!{sentence} फलित होता
  है (सार्वत्रिक निदर्शन, $x$ के लिए~$\num{m}$ और $y$ के
  लिए~$\num{n}$):}
& (\Obj Q_{q}(\num{m}, \num{n}) \land \Obj S_{\sigma}(\num{m},
\num{n})) \lif {}\\
& \qquad (\Obj Q_{q'}(\num{m}', \num{n}') \land
\Obj S_{\sigma'}(\num{m}, \num{n}') \land !A(\num{m}, \num{n})).
\intertext{आगमन-परिकल्पना से $!T(M, w) \Entails !C(M, w, n)$, अर्थात्}
& \Obj Q_q(\num{m}, \num{n}) \land \Obj S_{\sigma_0}(\num{0}, \num{n})
\land \dots \land \Obj S_{\sigma_k}(\num{k}, \num{n}) \land \\
& \qquad\lforall[x][(\num{k} < x \lif \Obj S_\TMblank(x, \num{n}))]\\
\intertext{$n$ चरणों के बाद पट्टी के खाना~$m$ में~$\sigma$ है, इसलिए
  संबंधित संयोजक $ \Obj S_\sigma(\num{m}, \num{n})$ है और इससे फलित
  होता है:}
& \Obj Q_{q}(\num{m}, \num{n}) \land \Obj S_{\sigma}(\num{m},
   \num{n})
\intertext{अब हमें मिलता है}
& \Obj Q_{q'}(\num{m}', \num{n}') \land \Obj S_{\sigma'}(\num{m},
  \num{n}') \land {}\\
& \qquad\Obj S_{\sigma_0}(\num{0}, \num{n}') \land \dots \land
  \Obj S_{\sigma_k}(\num{k}, \num{n}') \land {}\\
& \qquad \lforall[x][(\num{k} < x \lif \Obj S_\TMblank(x, \num{n}'))]
\end{align*}
इस प्रकार: पहली पंक्ति पूर्ववर्ती सशर्त के परिणामांश से प्रत्यक्षतः,
मोडस पोनेन्स द्वारा मिलती है। बीच की पंक्ति का प्रत्येक संयोजक---जिसमें
$S_{\sigma_m}(\num{m},\num{n}')$ सम्मिलित नहीं है---$!C(M, w, n)$ के
संबंधित संयोजक और $!A(\num{m}, \num{n})$ से मिलता है।

यदि $m < k$, तो $!T(M,w) \Proves \num{m} < \num {k}$
(\olref[rep]{prop:mlessk}), और $<$ की संक्रामकता से हमें
$\lforall[x][(\num{k} < x \lif \num{m} < x)]$ मिलता है। यदि $m = k$,
तो केवल तर्कशास्त्र से
$\lforall[x][(\num{k} < x \lif \num{m} < x)]$। तब अंतिम पंक्ति
$!C(M, w, n)$ के संबंधित संयोजक,
$\lforall[x][(\num{k} < x \lif \num{m} < x)]$, और
$!A(\num{m}, \num{n})$ से मिलती है। यदि $m<k$, तो यही पहले से
$!C(M, w, n+1)$ है।

अब मान लें कि $m=k$। उस दशा में $n+1$ चरणों के बाद पट्टी-शीर्ष ने
खाना~$k+1$ भी देख लिया है, जो अब देखा गया सबसे दायाँ खाना है। इसलिए
$!C(M, w, n+1)$ में एक नया संयोजक
$\Obj S_\TMblank(\num{k}',\num{n}')$ है, और अंतिम संयोजक
$\lforall[x][(\num{k}' < x \lif \Obj S_\TMblank(x, \num{n}'))]$ है।
हमें सत्यापित करना है कि ये दोनों !!{sentence} भी फलित होते हैं।

हमारे पास पहले से
$\lforall[x][(\num{k} < x \lif \Obj S_\TMblank(x, \num{n}'))]$ है।
विशेषतः इससे $\num{k} < \num{k}' \lif
\Obj S_\TMblank(\num{k}', \num{n}')$ मिलता है। अभिगृहीत
$\lforall[x][x < x']$ से हमें $\num{k} < \num{k}'$ मिलता है। मोडस
पोनेन्स से $\Obj S_\TMblank(\num{k}',\num{n}')$ फलित होता है।

साथ ही, चूँकि $!T(M,w) \Proves \num{k} < \num{k}'$, $<$ की
संक्रामकता का अभिगृहीत हमें
$\lforall[x][(\num{k}' < x \lif \Obj S_\TMblank(x, \num{n}'))]$ देता
है। (इसका सत्यापन अभ्यास के रूप में छोड़ा जाता है।)

\item मान लें कि \olref{left} रूप का कोई निर्देश है। तब
  \olref[rep]{defn:tm-descr}\olref[rep]{rep-left} के अनुसार,
\begin{align*}
& \lforall[x][\lforall[y][((\Obj Q_{q}(x', y) \land \Obj
    S_{\sigma}(x', y)) \lif {}]]\\
& \qquad (\Obj Q_{q'}(x, y') \land \Obj
  S_{\sigma'}(x', y') \land !A(x, y))) \land {}\\
& \lforall[y][((\Obj Q_{q_i}(\num{0}, y) \land \Obj S_{\sigma}(\num{0},
    y)) \lif {}]\\
& \qquad (\Obj Q_{q_j}(\num{0}, y') \land \Obj S_{\sigma'}(\num{0}, y')
  \land !A(\num{0}, y)))
\intertext{$!T(M,w)$ का एक संयोजक है। यदि $m>0$, तो $l = m - 1$
  (अर्थात् $m = l+1$) मानें। ऊपर के !!{sentence} के प्रथम संयोजक से
  निम्न फलित होता है:}
& (\Obj Q_{q}(\num{l}', \num{n}) \land \Obj S_{\sigma}(\num{l}', \num{n}))
\lif {} \\
& \qquad (\Obj Q_{q'}(\num{l}, \num{n}') \land \Obj S_{\sigma'}(\num{l}',
\num{n}') \land !A(\num{l}, \num{n}))
\intertext{अन्यथा $l = m = 0$ मानें और दूसरे संयोजक से फलित निम्न
  !!{sentence} पर विचार करें:}
& ((\Obj Q_{q_i}(\num{0}, \num{n}) \land \Obj S_{\sigma}(\num{0}, \num{n})) \lif {}\\
& \qquad   (\Obj Q_{q_j}(\num{0}, \num{n}') \land \Obj S_{\sigma'}(\num{0}, \num{n}') \land
{}!A(\num{0}, \num{n})))
\intertext{दोनों में से कोई भी वाक्य निम्न को फलित करता है:}
&  \Obj Q_{q'}(\num{l}, \num{n}') \land \Obj S_{\sigma'}(\num{m},
  \num{n}') \land {}\\
&\qquad  \Obj S_{\sigma_0}(\num{0}, \num{n}')\land \dots \land
  \Obj S_{\sigma_k}(\num{k}, \num{n}')  \land {} \\
& \qquad  \lforall[x][(\num{k} < x
  \lif \Obj S_\TMblank(x, \num{n}'))]
\end{align*}
पहले की तरह। (ध्यान दें कि पहली दशा में
$\num{l}' \ident \num{l+1} \ident \num{m}$, और दूसरी दशा में
$\num{l} \ident \num{0}$।) किंतु यही $!C(M, w, n+1)$ है।

\item स्थिति \olref{stay} को अभ्यास के रूप में छोड़ा जाता है।
\end{enumerate}
हमने दिखा दिया कि किसी भी~$n$ के लिए
$!T(M, w) \Entails !C(M, w, n)$।
\end{proof}

\begin{prob}
\olref[tur][und][ver]{lem:config} के प्रमाण की
स्थिति~\olref[tur][und][ver]{stay} पूरी करें।
\end{prob}

\begin{prob}
$\Obj S_{\sigma_i}(\num{i}, \num{n})$ और $!A(m, n)$ से
$\Obj S_{\sigma_i}(\num{i}, \num{n}')$ की !!a{derivation} दें (मानते
हुए कि $i \neq m$, अर्थात् या तो $i < m$ अथवा $m < i$)।
\end{prob}

\begin{prob}
$\lforall[x][(\num{k} < x \lif \Obj S_\TMblank(x, \num{n}'))]$,
$\lforall[x][x < x']$, और
$\lforall[x][\lforall[y][\lforall[z][ ((x < y \land y < z) \lif x <
      z)]]]$ से
$\lforall[x][(\num{k}' < x \lif \Obj S_\TMblank(x, \num{n}'))]$ की
!!a{derivation} दें।)
\end{prob}


\begin{lem}
\ollabel{lem:valid-if-halt}
यदि $M$ आगत~$w$ पर रुकती है, तो $!T(M, w) \lif !E(M, w)$ वैध है।
\end{lem}

\begin{proof}
\olref{lem:config} से हम जानते हैं कि किसी भी समय~$n$ के लिए समय~$n$
पर $M$ के विन्यास का वर्णन~$!C(M, w, n)$, $!T(M, w)$ से फलित होता है।
मान लें कि $M$, $k$ चरणों के बाद रुकती है। उस समय वह किसी
$m \in \Nat$ के लिए खाना~$m$ पढ़ रही होगी। तब $!C(M, w, k)$, $M$ के
रुकने वाले विन्यास का वर्णन करता है, अर्थात् इसमें
$\Obj Q_q(\num{m}, \num{k})$ और $\Obj S_\sigma(\num{m}, \num{k})$
दोनों संयोजक होते हैं, जहाँ $\delta(q,\sigma)$ अपरिभाषित है। अतः
\olref{lem:halt-config-implies-halt} के अनुसार
$!C(M, w, k) \Entails !E(M, w)$। किंतु चूँकि
$!T(M, w) \Entails !C(M, w, k)$, हमारे पास
$!T(M, w) \Entails !E(M, w)$ है, और इसलिए
$!T(M, w) \lif !E(M, w)$ वैध है।
\end{proof}


\begin{explain}
अपने दावे का सत्यापन पूरा करने के लिए हमें विपरीत दिशा भी स्थापित करनी
है: यदि $!T(M, w) \lif !E(M, w)$ वैध है, तो आगत~$w$ पर आरंभ किए जाने
पर $M$ वास्तव में रुकती है।
\end{explain}

\begin{lem}
\ollabel{lem:halt-if-valid}
यदि $\Entails !T(M, w) \lif !E(M, w)$, तो $M$ आगत~$w$ पर रुकती है।
\end{lem}

\begin{proof}
प्रांत~$\Nat$ वाली $\Lang L_M$-!!{structure}~$\Struct M$ पर विचार
करें, जो $\Obj 0$ की व्याख्या~$0$ के रूप में, $\prime$ की उत्तरवर्ती
फलन के रूप में, और $<$ की लघुतर संबंध के रूप में करती है; और
विधेयों $\Obj Q_q$ तथा~$\Obj S_\sigma$ की व्याख्या इस प्रकार करती
है:
\begin{align*}
  \Assign{\Obj Q_q}{M} & =
\Setabs{\tuple{m, n}}{\begin{array}{ll}\text{$w$ पर आरंभ करने के बाद, $n$ चरणों पश्चात्,}\\ \text{$M$ अवस्था $q$ में है
  और खाना~$m$ पढ़ रही है}\end{array}} \\
\Assign{\Obj S_\sigma}{M} & = \Setabs{\tuple{m, n}}{\begin{array}{ll}
\text{$w$ पर आरंभ करने के बाद, $n$ चरणों पश्चात्,}\\ \text{$M$ के खाना~$m$ में
  प्रतीक~$\sigma$ है}\end{array}}
\end{align*}
दूसरे शब्दों में, हम !!{structure}~$\Struct{M}$ की रचना इस तरह करते हैं
कि वह आगत~$w$ पर आरंभ की गई $M$ के वास्तविक व्यवहार को चरण-दर-चरण
वर्णित करे। स्पष्ट है कि $\Sat{M}{!T(M, w)}$। यदि
$\Entails !T(M, w) \lif !E(M, w)$, तो $\Sat{M}{!E(M, w)}$ भी, अर्थात्
\[
\Sat{M}{\lexists[x][\lexists[y][(\bigvee_{\tuple{q, \sigma} \in
      X}(\Obj Q_q(x, y) \land \Obj S_\sigma(x, y)))]]}.
\]
चूँकि $\Domain{M} = \Nat$, कोई $m$, $n \in \Nat$ ऐसे अवश्य हैं कि
किसी ऐसे~$q$ और~$\sigma$ के लिए
$\Sat{M}{\Obj Q_q(\num{m}, \num{n}) \land \Obj S_\sigma(\num{m},
\num{n})}$, जिनके लिए $\delta(q, \sigma)$ अपरिभाषित है।
$\Struct M$ की परिभाषा के अनुसार इसका अर्थ है कि आगत~$w$ पर आरंभ की
गई $M$, $n$ चरणों के बाद अवस्था~$q$ में है और प्रतीक~$\sigma$ पढ़ रही
है, तथा संक्रमण फलन अपरिभाषित है; अर्थात् $M$ रुक चुकी है।
\end{proof}

\end{document}
